\Hefteintrag{2}{Programmablauf grafisch darstellen: Struktogramme}{%
Struktorgramme (auch: \LoesungLuecke{Nassi-Shneiderman-Diagramme}{12cm}) stellen \emphColA{Algorithmen (=Ablauf eines Programms)} grafisch dar. Ein Struktogramm bezieht sich immer auf genau eine Methode und stellt aufgerufene Methoden als ein Element dar.

In \textbf{Hefteintrag \ref{sec:befehlsarten}} sind die Struktogramm-Bausteine für verschiedene Befehlsarten dargestellt, die beliebig miteinander kombiniert werden können. 

Zusätzlich gibt es folgenden Baustein, der eine leere Stelle (z.B. leerer else-Block) darstellt:

\vspace{4cm}


}